% ==============================================================================
% Secção 2: Descrição sumária da proposta
% ==============================================================================

\subsection{Resumo do projeto (em português)}
\begin{campotexto}
Inserido no horizonte da engenharia daqui a 20 anos, este projeto tem como objetivo conceber soluções para desenvolver e otimizar os carros do futuro, respondendo aos desafios de eficiência energética, segurança, conforto e sustentabilidade da mobilidade elétrica. Partindo das limitações das baterias de lítio, das suspensões passivas e da recuperação de energia restrita à travagem, o trabalho segue cinco etapas — questões orientadoras, revisão da literatura, análise PESTEL, \textit{benchmarking} e síntese crítica — para propor uma abordagem integrada. No armazenamento de energia, exploram-se baterias de estado sólido e de sódio e respetivas opções de descarte. Para a dinâmica do veículo, projeta-se a integração de suspensões ativas, sensores com inteligência artificial para atuação preditiva e amortecedores regenerativos que convertem vibrações em eletricidade. Conclui-se que o carro do futuro exigirá a conjugação destas tecnologias e a validação da sua viabilidade industrial.
\end{campotexto}

\subsection{Abstract}
\begin{campotexto}
Set within the horizon of engineering 20 years from now, this project aims to develop and optimize the cars of the future, addressing the energy efficiency, safety, comfort, and sustainability challenges of electric mobility. Moving beyond the limitations of dominant lithium batteries, passive suspensions, and braking-only energy recovery, the work follows five stages — guiding questions, literature review, PESTEL societal analysis, benchmarking, and critical synthesis — to propose an integrated technological pathway. For energy storage, it explores solid-state and sodium batteries alongside their disposal options. For vehicle dynamics, it envisions integrating active suspensions, sensor- and AI-based predictive control, and regenerative shock absorbers that convert vibrations into electricity. Ultimately, building the vehicle of the future will require uniting these emerging technologies and validating their industrial viability and impact on range.
\end{campotexto}

\subsection{Resumo Visual (Graphical Abstract)}
\begin{campotexto}
\instrucao{Graphical abstract in a single image or diagram that visually summarizes your project, considering concepts, technological innovations, methodologies and conclusions, providing a quick, at-a-glance overview for readers. Its purpose is to enhance online browsing and help evaluators and investors to quickly identify relevant advantages by presenting the core message in a clear and concise, self-explanatory format.}
\end{campotexto}
